\documentclass[a4paper]{article}

\usepackage{listings}
\usepackage[margin=1in]{geometry} % full-width
\usepackage{biblatex}
\usepackage{graphicx}

% setup for code highlighting, from Google search + Gemini
\usepackage{xcolor}
\definecolor{commentsColor}{rgb}{0.497495, 0.497587, 0.497464}
\definecolor{keywordsColor}{rgb}{0.000000, 0.000000, 0.635294}
\definecolor{stringColor}{rgb}{0.558215, 0.000000, 0.135316}

\lstset{
    language=Python,
    numbers=left,
    numberstyle=\tiny\color{gray},
    keywordstyle=\color{keywordsColor}\bfseries,
    commentstyle=\color{commentsColor}\itshape,
    stringstyle=\color{stringColor},
    tabsize=2,
    title=\lstname,
    columns=fixed
}


\addbibresource{refs.bib}


\title{Statistical Efficiency of Birth-Death Tree Parameter Estimation}
\author{Santosh Desai \and Matt McAnear \and Urvi Mehta}

\date{
	\today
}

\begin{document}
	\maketitle

Source code available at:

PhylodynamicsDL: https://github.com/mcanearm/PhylodynamicsDL

PhyloDeepPOMP: https://github.com/Horopter/PhyloDeepPOMP

\section{Intro}

The Phylodeep package implements a deep learning approach to estimating parameters of several different
tree models. The results of the deep learning models are impressive, but noticeably absent from 
the baseline of the models is the closed-form MLE for any of the simulated trees. For our study, we focused on 
the Birth-Death trees and compared the statistical efficiency of
the Phylodeep approach to the traditional Maximum Likelihood Estimation (MLE) approach. Initial comparisons 
indicated that there may be some issues with the likelihood estimation component, but we think the results
are interesting enough to present, albeit with caveats.

\section{Data Generation}

The original paper by Voznicka et. al.\cite{voznica_deep_2022} utilized a custom simulator to generate trees using a package by Zhukova et. al.\cite{zhukova_evolbioinfotreesimulator_2025},
one of the other authors. We chose to replicate this process as faithfully as possible, both for 
reasons of simplicity and reproducibility.

\begin{lstlisting}[language=Python, caption=Simulator Code]
subprocess.run(
    [
        str(generate_bd), "--min_tips", "10", "--max_tips", "500",
            "--la", str(la), "--psi", str(psi), "--p", str(0.5), 
            "--nwk", tree_path, "--log", log_path,
    ],
    check=True,
    timeout=timeout_seconds,
    stdin=subprocess.DEVNULL,
    stdout=subprocess.DEVNULL,
    stderr=subprocess.DEVNULL,
)
\end{lstlisting}

This is identical to the process outlined in the Phylodeep repo\cite{zhukova_phylodeepsimulators_nodate},
with the exception that we reduce the minimum tips to 10. Our random sampling of $\lambda$ and $\psi$ are between
0 and 1. These ranges do not align with the Phlyodeep paper in the supplementary data, and
we locked the sampling probability at 0.5 rather than sampling it as well. By focusing on 
two parameters, we are simplifying our analysis further, which seemed reasonable given
the timeframe of the project.

\section{Methods}

\section{MLE}

We made a few attemps at implementing the MLE, which ended up being a more challenging part of the project
than we expected, mostly due to difficulties in dealing with the data. 

URVI - put in what you did here and results / final outcome


\section{Phylodeep}

SANTOSH - model fitting process

\section{Results}

Using the MLE and Phylodeep model together, we found that the Phylodeep model consistently outperformed
the MLE method in estimating the main parameters ($mu$ and $lambda$) of the Birth-Death trees. Surprisingly,
the overall tree size does seem to reduce the RMSE for both Phylodeep and MLE methods, but 
not by much.

\begin{figure}
\centering
\includegraphics[width=0.9\linewidth]{images/RMSE_vs_tips.png}
\caption{Bootstrap estimated RMSE of $\lambda$ and $\mu$ estimates from MLE and Phylodeep models across varying tree sizes.}
\end{figure}

\begin{figure}
\centering
\includegraphics[width=0.9\linewidth]{images/estimation_error_vs_tips.png}
\caption{Raw error estimates for $\lambda$ and $\mu$ from MLE and Phylodeep models. There appears to be no relationship with the size of the tree in this context, but the MLE is consistently biased upwards.}
\label{fig:bias_plot}
\end{figure}

\begin{figure}
    \centering
    \includegraphics[width=0.8\linewidth]{images/lambda_vs_mu_estimates.png}
    \caption{Scatter plot of $\lambda$ vs $\mu$ estimates from MLE and Phylodeep models. The MLE estimates land on the lower boundary condition for $\mu$ close to 0, indicating an issue of identifiability with our current implementation.}
\end{figure}

When we consider the fitted values of the Phylodeep estimation vs. the MLE estimation, we see
a pattern. Plotting the two parameter estimates together, we see that the MLE estimates often
seem to settle on a boundary condition for $mu$, while the Phylodeep estimates
do not have this issue. Further, Phylodeep estimates tend to miss equally both above and 
below the true parameter, whereas the MLE values are notably positive (see Figure \ref{fig:bias_plot}).


\section{Conclusion}

Our main research question was around the statistical efficiency of these two competing estimators
with regards to data size (number of tips in the tree) in the context of the Birth-Death tree model.
In our initial attempt at answering this question, it appears that Phylodeep outperforms the MLE in 
most contexts for all tree sizes.

\break

\printbibliography
	
\end{document}